% Original teaching diagram, 2026-11-05. No background grid.
\begin{tikzpicture}[x=1mm,y=1mm,>=Stealth,font=\figurefont\fontsize{8}{10}\selectfont,
 v/.style={circle,draw=FigureStroke,fill=white,minimum size=6mm,inner sep=0pt,line width=.75pt},
 bag/.style={rounded corners=1.8pt,draw=FigureBlue,fill=FigureBlueFill,minimum width=24mm,minimum height=10mm,inner sep=2pt},
 e/.style={draw=FigureStroke,line width=.75pt},
 ar/.style={->,draw=FigureBlue,line width=.9pt},
 lab/.style={text=FigureMuted,font=\figurefont\fontsize{7.5}{9}\selectfont}]

\fill[FigureBlueFill,rounded corners=2pt] (2,-5) rectangle (47,31);
\node[v,fill=FigureYellowFill] (c) at (10,24) {丙};
\node[v] (a) at (10,2) {甲};\node[v] (one) at (39,13) {$1$};
\node[v] (b) at (69,24) {乙};\node[v] (two) at (97,24) {$2$};\node[v,fill=FigureYellowFill] (three) at (97,2) {$3$};
\draw[ar,dashed,draw=FigureGreen] (c)--(one);
\draw[draw=FigureBlue,line width=1.4pt] (a)--(one);
\draw[->,draw=FigureGreen,line width=.9pt] (one) to[bend left=18] (a);
\draw[draw=FigureBlue,line width=1.4pt] (b)--(two);\draw[e] (b)--(three);
\node[lab,above] at (25,31) {$Z$};

\end{tikzpicture}
